% ===================================================================
% Section 3: PI regulátor 
% ===================================================================

\section{PI regulátor pre systém TS05}
\label{sec:pi_controller}

V tejto kapitole je riešený návrh a implementácia PI regulátora
ako základného nástroja pre reguláciu tepelnej sústavy.
Najskôr je stručne zhrnutá teoretická štruktúra PI regulátora
a jeho zapojenie v uzavretom regulačnom obvode,
následne je tento prístup aplikovaný na systém TS05
vrátane ladenia parametrov a realizácie v reálnom čase.

% -------------------------------------------------------------------
% Subsection 3.1 - Štruktúra PI regulátora
% -------------------------------------------------------------------
\subsection{Štruktúra PI regulátora}

PI regulátor predstavuje zjednodušenú formu PID regulátora,
v ktorej je derivačná zložka vynechaná.
Pre systémy s pomalou dynamikou, ako sú tepelné sústavy,
derivačný člen spravidla neprináša zlepšenie regulácie
a naopak môže zosilňovať merací šum.

Paralelná forma PI regulátora je definovaná prenosovou funkciou
\begin{equation}
	C(s) = K_p + \frac{K_i}{s},
	\label{eq:pi_parallel}
\end{equation}
kde $K_p$ je proporcionálny zisk a $K_i$ integračný zisk regulátora.

Regulátor je zapojený v štandardnej uzavretej spätnej väzbe
so zápornou spätnou väzbou.
Bloková schéma uzavretého regulačného obvodu je znázornená
na obr.~\ref{fig:pi_block_scheme}.

\begin{center}
	\vbox{
	\makebox[\textwidth][c]{%
		\includegraphics[width=1.2\textwidth]{TS_uro_block_scheme.pdf}
	}
	\figcaption{Bloková schéma uzavretého regulačného obvodu s PI regulátorom}
	\label{fig:pi_block_scheme}
	}
\end{center}

\paragraph{Prenosy regulačného obvodu}

Nech $G_s(s)$ označuje prenos regulovaného procesu
a $C(s)$ prenos regulátora.
Prenos otvoreného regulačného obvodu je potom daný vzťahom
\begin{equation}
	G_{\mathrm{ORO}}(s) = C(s)\,G_s(s),
	\label{eq:pi_oro_general}
\end{equation}
pričom prenos uzavretého regulačného obvodu má tvar
\begin{equation}
	G_{\mathrm{URO}}(s) =
	\frac{C(s)\,G_s(s)}{1 + C(s)\,G_s(s)}.
	\label{eq:pi_uro_general}
\end{equation}

Stabilita uzavretého regulačného obvodu je daná koreňmi
charakteristickej rovnice
\begin{equation}
	1 + C(s)\,G_s(s) = 0.
	\label{eq:pi_char_eq}
\end{equation}
Pre spojité systémy je podmienkou stability,
aby všetky póly uzavretého obvodu mali zápornú reálnu časť.

\paragraph{Zákon riadenia}

V časovej oblasti možno PI regulátor zapísať v tvare
\begin{equation}
	u(t) = K_p\,e(t) + K_i \int_0^t e(\tau)\,d\tau,
	\label{eq:pi_control_law}
\end{equation}
kde $e(t) = w(t) - y(t)$ predstavuje regulačnú odchýlku,
$w(t)$ žiadanú hodnotu a $y(t)$ výstup systému.

Pre diskrétnu realizáciu s periódou vzorkovania $T_s$
je integračný člen aproximovaný rekurentným vzťahom
\begin{align}
	e(k) &= w(k) - y(k), \\
	I(k) &= I(k-1) + T_s\,e(k), \\
	u(k) &= K_p\,e(k) + K_i\,I(k).
\end{align}

Z dôvodu fyzikálnych obmedzení akčného zásahu
je v praxi uplatnená saturácia riadiaceho signálu.

\paragraph{Pseudokód PI regulátora}

\begin{lstlisting}[language=Matlab]
% Inputs: w (setpoint), y (output), Ts
% States: I (integrator state)
% Params: Kp, Ki, u_min, u_max

e = w - y;

% Integrator update
I = I + Ts * e;

% PID law
u = Kp * e + Ki * I;

% Saturation
u = min(max(u, u_min), u_max);

\end{lstlisting}

% -------------------------------------------------------------------
% Subsubsection 3.2 Ladenie PI regulátora
% -------------------------------------------------------------------

\subsection{Nastavenie parametrov PI regulátora pre systém TS05}

Pre získanie parametrov regulátora bola využitá funkcia \texttt{pidtune} v prostredí MATLAB, ktorá vykonáva automatizované ladenie regulátora.  
Tento postup bol zvolený preto, že viaceré analytické metódy ladenia nie je možné na systém TS05 aplikovať.

Charakteristický polynóm uzavretého regulačného obvodu pre PI regulátor má po dosadení tvar

\begin{equation}
	s^3 
	+ \frac{b_1}{b_2} s^2 
	+ \left( \frac{b_0}{b_2} + \frac{a_0 K_p}{b_2} \right) s 
	+ \frac{a_0 K_i}{b_2} 
	= 0
	\label{eq:cl_poly_real}
\end{equation}

Ak by sme chceli použiť metódu umiestnenia pólov pomocou rozšíreného želaného polynómu tretieho rádu, jeho všeobecný tvar je

\begin{equation}
	s^3 + (p + 2 k \omega_0) s^2 + (2 k \omega_0 p + \omega_0^2) s + \omega_0^2 p = 0
	\label{eq:desired_poly}
\end{equation}

Porovnaním koeficientov polynómov \eqref{eq:cl_poly_real} a \eqref{eq:desired_poly} dostaneme sústavu rovníc:

\begin{align}
	1 &= 1
	\label{eq:compare1} \\
	p + 2 k \omega_0 &= \frac{b_1}{b_2}
	\label{eq:compare2} \\
	2 k \omega_0 p + \omega_0^2 &= \frac{b_0}{b_2} + \frac{a_0 K_p}{b_2}
	\label{eq:compare3} \\
	\omega_0^2 p &= \frac{a_0 K_i}{b_2}
	\label{eq:compare4}
\end{align}

Táto sústava nemá realizovateľné riešenie pre parametre $K_p$ a $K_i$.  
Dôvodom je, že členy závislé od $K_p$ a $K_i$ sa vyskytujú iba v dvoch rovniciach, zatiaľ čo parametre želaného dynamického správania $k$, $\omega_0$, $p$ ovplyvňujú tri rovnice.  
Sústava je preto pre fyzikálne konzistentné riešenie preurčená a neexistuje v nej žiadna kombinácia parametrov, ktorá by spĺňala všetky rovnice súčasne.

Rovnako nie je možné použiť metódu Ziegler–Nichols, keďže námi identifikovaný model (\ref{eq:Gs_ls_approx}) nemá kritický bod, ktorý je pre túto metódu potrebný.

Keďže návrh PI regulátora nie je hlavnou témou práce, bol zvolený numerický nástroj \texttt{pidtune}.  
Jeho princíp spočíva v tom, že pre zvolenú požadovanú šírku pásma uzavretého systému vypočíta optimálne hodnoty parametrov PI regulátora tak, aby bola zabezpečená dobrá dynamická odozva a primeraná robustnosť.

Najskôr bol systém prevedený do diskrétnej podoby:

\begin{lstlisting}[language=Matlab]
Gs = Gs_fmin;
Ts = 0.1;

Gz = c2d(Gs, Ts, 'tustin');
[Ad,Bd,Cd,Dd] = ssdata(Gz);
\end{lstlisting}

Následne bolo vykonané ladenie PI regulátora:

\begin{lstlisting}[language=Matlab]
wc = 0.5;
Cz = pidtune(Gz, 'pi', wc);

[numC,denC] = tfdata(Cz,'v');
[Ac,Bc,Cc,Dc] = tf2ss(numC,denC);
\end{lstlisting}

Získané parametre regulátora:

\[
K_p = 3.58, \qquad K_i = 1.84
\]

Prenosová funkcia otvoreného regulačného obvodu je

\begin{equation}
	G_{oro}(s) =
	\frac{21.21 s + 11.15}{s^3 + 59.01 s^2 + 8.95 s}
	\label{eq:oro_tf}
\end{equation}

Póly otvoreného systému sú:

\begin{equation}
	s_{1} = 0, \quad s_{2} = -58.855, \quad s_{3} = -0.1521
	\label{eq:oro_poles}
\end{equation}

Znakom stability je prítomnosť pólov s negatívnou reálnou časťou, čo je splnené.

Simulácia bola vykonaná s obmedzením akčného zásahu na interval \([-4, 6]\),  
čo zodpovedá pracovnému bodu 4 V.  
Výsledok simulácie je uvedený na obrázku~\ref{fig:pi_simulation_results}.

\begin{center}
	\vbox{
		\makebox[\textwidth][c]{%
			\includegraphics[width=1.5\textwidth]{ts_pb_vent6_spir4_pidtune_sim.png}
		}
		\figcaption{Simulačný výsledok uzavretého PI riadenia systému TS05}
		\label{fig:pi_simulation_results}
	}
\end{center}

Z výsledkov boli namerané tieto hodnoty:

\begin{align}
	\text{Čas ustálenia} &= 11.5 \text{ s}
	\label{eq:settling_time} \\
	\text{Prekmit} &= 15.9 \%
	\label{eq:overshoot} \\
	\text{MSE} &= 0.04369
	\label{eq:mse_pi}
\end{align}

Dosiahnutý výsledok možno považovať za vhodný pre riadenie tepelného systému.  
Cieľom je dosiahnuť čo najrýchlejšiu odozvu.  
Prekmit približne 16 percent je akceptovateľný, keďže riadený skok je malý a tepelná sústava je pomalá, čo prirodzene limituje tvar dynamickej odozvy.

% -------------------------------------------------------------------
% Subsubsection 3.3 Realizácia PI regulátora pre riadenie v reálnom čase
% -------------------------------------------------------------------

\subsection{Implementácia PI regulátora v MATLABe}

Ako je uvedené v zadaní, cieľom práce je implementácia riadiacich algoritmov na rôznych hardvérových a softvérových platformách.  
Preto bol PI regulátor implementovaný tak, aby mohol pracovať v reálnom čase bez použitia hotových Simulink blokov určených na riadenie.  
Celá logika riadenia bola zapísaná priamo do bloku MATLAB Function, ktorý sa vykonáva raz za každý riadiaci cyklus s periódou vzorkovania 0.1 sekundy.

Implementácia regulátora je nasledovná:

\begin{lstlisting}[language=Matlab]
function [u, integrator_out, y_ss_out] = ...
         pi_block(y, setpoint, integrator_in, t, y_ss_in)

% PI controller

pb_u = 4;
Ts = 0.1;
u_min = - pb_u;
u_max = 10 - pb_u;

Kp = 3.58;
Ki = 1.84;

% waiting phase for steady-state detection
if t < 10*60
    u = 4;
    y_ss_out = y;
    integrator_out = 0;

else
    y_ss_out = y_ss_in;
    pb_y = y_ss_in;

    % output normalization
    y = y - pb_y;

    % control error
    e = setpoint - y;

    % integrator
    integrator_out = integrator_in + e * Ts;

    % PI control law
    u = Kp * e + Ki * integrator_out;
    u = min(max(u, u_min), u_max);

    % input normalization
    u = u + pb_u;

end
end
\end{lstlisting}
  
Tento blok však na rozdiel od štandardných Simulinkových riadiacich blokov neobsahuje vlastnú vnútornú pamäť, preto je potrebné stavové premenné prenášať medzi jednotlivými volaniami bloku explicitne.

Z tohto dôvodu sú hodnoty integrátora a odhad ustáleného stavu odosielané na výstup PI bloku a v nasledujúcom kroku sú privádzané späť na vstup.  
Týmto spôsobom sa zabezpečí, že regulátor si zachováva svoje vnútorné stavy z jedného riadiaceho kroku do druhého.

Logika zapojenia bloku MATLAB Function je zobrazená na obrázku~\ref{fig:pi_rt_scheme},  
ktorý ukazuje vstupy, výstupy a spätnú väzbu stavových premenných potrebných pre správnu činnosť regulátora.

\begin{center}
    \vbox{
        \makebox[\textwidth][c]{%
            \includegraphics[width=1.3\textwidth, trim=100 150 100 150, clip]{TS_pi_scheme.pdf}
        }
        \figcaption{Zapojenie PI regulátora v bloku MATLAB Function so spätnou väzbou stavov}
        \label{fig:pi_rt_scheme}
    }
\end{center}

Pred samotným začiatkom regulácie bol zavedený časový úsek, počas ktorého regulátor nevykonáva žiadny zásah.  
Počas prvých desiatich minút sa na výstupe drží hodnota pracovného bodu a zároveň sa sleduje teplota na výstupe systému.  
Cieľom je získať reálnu hodnotu ustáleného stavu priamo z aktuálnych podmienok experimentu.

Použitie hodnôt zosadených zo statických charakteristík (obr.~\ref{fig:ts_static_sensors6}) nie je vhodné.  
Tepelný systém je citlivý na teplotu v miestnosti a experimenty nebolo možné realizovať v rovnakých podmienkach.  
Okrem toho surové merania (obr.~\ref{fig:ts_pb_vent6_spir4_ident_raw}) ukázali, že teplota narastá veľmi pomaly aj počas dlhého experimentu, čo znamená, že systém nemá pevnú, opakovateľnú hodnotu ustáleného stavu.

Z toho dôvodu je ustálený stav zisťovaný automaticky priamo regulátorom ako posledná nameraná hodnota $y$ pred spustením PI riadenia.  
Tento prístup funguje bez ohľadu na to, či je systém iba málo zahriaty alebo už dosiahol vysokú prevádzkovú teplotu.

\paragraph{Experiment 1: Regulácia okolo jedného pracovného bodu}

Na základe implementácie uvedenej vyššie bol následne vykonaný experiment v simulačnom prostredí s cieľom overiť kvalitu regulácie.  
Priebeh regulácie je zobrazený na obrázku~\ref{fig:ts_pi_real_basic}.

\begin{center}
    \vbox{
        \makebox[\textwidth][c]{%
            \includegraphics[width=1.5\textwidth]{ts_pb_vent6_spir4_PI_real.png}
        }
        \figcaption{Simulačný experiment s PI regulátorom pri skokovej zmene žiadanej hodnoty}
        \label{fig:ts_pi_real_basic}
    }
\end{center}

Kvalita regulácie bola hodnotená podľa troch kritérií:

\begin{itemize}
    \item prekmit: žiadaná hodnota bola 0.35~V, maximálna hodnota výstupu dosiahla približne 0.52~V,  
    čo zodpovedá približne 49.0 percentnému prekmitu,
    \item čas ustálenia: výstup sa ustálil približne po 28 sekundách,
    \item trvalá regulačná odchýlka: prakticky nulová.
\end{itemize}

Pri porovnaní so simuláciou uvedenou v predchádzajúcich kapitolách možno vidieť, že prekmit je výrazne vyšší.  
Dôvodom je skutočnosť, že identifikované modely nedokázali zachytiť celú dynamiku systému TS05 a viacero pomalých javov,  
ktoré sa prejavujú až pri reálnej prevádzke.  
Napriek tomu je regulačný čas veľmi priaznivý vzhľadom na to, že ide o pomalý tepelný systém.  
Rýchla regulácia bola prioritou, preto možno výsledok považovať za uspokojivý pre riadenie okolo jedného pracovného bodu.

\paragraph{Experiment 2: Reakcia na väčší skok}

Nasledoval druhý experiment, ktorého cieľom bolo overiť správanie regulátora pri väčšej skokovej zmene žiadanej hodnoty, 
ktorá má simulovať prechod medzi rôznymi pracovnými bodmi systému.  
Systém bol najprv ustálený v pracovnom bode s výstupnou hodnotou približne 6.4~V a následne bola žiadaná hodnota skokovo zmenená na 7.5~V. 
Výsledok je zobrazený na obrázku~\ref{fig:ts_pi_real_big_step}.

\begin{center}
    \vbox{
        \makebox[\textwidth][c]{%
            \includegraphics[width=1.5\textwidth]{ts_pb_vent6_pbTopb_real_badToShow.png}
        }
        \figcaption{Experiment PI regulátora pri pri skoku z 6.4~V na 7.5~V}
        \label{fig:ts_pi_real_big_step}
    }
\end{center}

Z grafu je zrejmé, že počas regulácie došlo k výraznému presýteniu akčného zásahu.  
Ovládací signál dosiahol svoj limit a počas určitého času na ňom zotrval.  
Aj keď bola požadovaná hodnota dosiahnutá, takéto správanie je z pohľadu bezpečnosti a kvality regulácie nevyhovujúce.

Z tohto dôvodu možno konštatovať, že navrhnutý PI regulátor je vhodný na riadenie v blízkosti jedného pracovného bodu,  
avšak nie je vhodný na prepínanie medzi vzdialenejšími pracovnými bodmi.  
V ďalšej časti práce bude preto implementovaný postup umožňujúci bezpečné a stabilné prepínanie medzi pracovnými bodmi  
a následne budú jednotlivé regulátory skombinované tak, aby systém disponoval širším využiteľným rozsahom riadenia.